% Near-root full-corner proof and disjoint range assembly for Section 7.
% This replaces SECTION7_FULL_CORNER_V3 inside the proof of Lemma 7.1.

\displayheading{Full corner}

Put $h=k-\ell$ and write
\[
  v(h):=\sum_i u_i h_i=n-m.
\]
The full corner is the range
\begin{equation*}
  v(h)\le\frac n{32}.
  \tag{7.26a}
\end{equation*}
We use the exact reverse recurrence (7.6) without replacing any ambient
falling factorial.

First make the uniform smallness of the residual first moments explicit. For
any $0\le w\le n$ and every one of the four block sizes,
\[
  \frac{\mu_{u_i}(w)}{\mu_{u_i}(n)}
  =\frac{(w)_{u_i}}{(n)_{u_i}}
  \le\left(\frac wn\right)^{u_i}.
\]
Equations (2.2) and (2.8) give
$\mu_{u_i}(n)\le n^{6+o(1)}$ uniformly in the phase. Moreover,
$u_i=(2/\log 2+o(1))\log n$, and
$\log 32=5\log 2$. Consequently, uniformly for $w\le n/32$,
\begin{equation*}
  \mu_{u_i}(w)
  \le
  n^{6+o(1)}32^{-u_i}
  =n^{-4+o(1)}.
  \tag{7.26}
\end{equation*}
In particular, there is one phase-independent threshold after which
\begin{equation*}
  \mu_{u_i}(w)\le n^{-3}
  \qquad
  (w\le n/32,\ 2\le i\le5).
  \tag{7.26b}
\end{equation*}

Fix a final residual profile $h$ satisfying (7.26a), and expose its blocks in
any order. Let $a$ be the current residual profile and
$v(a)=\sum_i u_i a_i$. If the next block has type $i$, then
$a_i<h_i$ and therefore
\[
  v(a)+u_i\le v(h)\le n/32.
\]
The exact recurrence (7.6), $k_i-a_i\le k_i\le n$, and (7.26b) give
\[
  \frac{B(a+e_i)}{B(a)}
  =
  \frac{2(k_i-a_i)\mu_{u_i}(v(a)+u_i)}{(a_i+1)^2}
  \le2n^{-2}<1
\]
for all sufficiently large $n$. Hence $B$ never increases along the path from
$0$ to $h$, and
\begin{equation*}
  B(h)\le B(0)=1.
  \tag{7.27a}
\end{equation*}
By the exact complementary identity (7.5) and the one-part-buffer first-moment
margin (5.19),
\begin{equation*}
  D(k-h)
  =\frac{B(h)}{\mathbb E Z_{\mathbf k}^{\mathrm{sgn}}}
  \le
  \exp\{-c_{\mathrm{nr}}(\log n)^2\}.
  \tag{7.27}
\end{equation*}
Equivalently, each residual profile contributes at most
$e^{-c_{\mathrm{nr}}(\log n)^2}$.
There are at most $(k_{\mathrm{co}}+1)^4$ residual profiles. Therefore
\[
  \sum_{v(h)\le n/32}D(k-h)
  \le
  \exp\!\left\{
    4\log(k_{\mathrm{co}}+1)-c_{\mathrm{nr}}(\log n)^2
  \right\}
  =o(1)
\]
uniformly in the phase. The quadratic logarithmic margin is stronger than the
$O(\log n)$ logarithm of the profile count. Equivalently, the full-corner sum
is bounded by one deterministic sequence
$\varepsilon_n^{\mathrm{full}}\to0$.

\displayheading{Disjoint three-range assembly}

For sufficiently large $n$, one has $\eta<31/32$. The coordinate box of all
partial subprofiles is then the disjoint union of
\[
\begin{aligned}
  \mathcal E_n&:=\{\ell:m\le\eta n\},\\
  \mathcal C_n&:=\{\ell:m>\eta n,\ n-m>n/32\},\\
  \mathcal F_n&:=\{\ell:n-m\le n/32\}.
\end{aligned}
\]
Indeed, if $\ell\in\mathcal F_n$, then
$m\ge31n/32>\eta n$, so the full corner cannot meet the empty corner; after
excluding those two cases, the two strict inequalities defining
$\mathcal C_n$ are automatic. Thus no boundary term is counted twice and no
subprofile is omitted.

The three estimates proved above have the form
\[
\begin{aligned}
  1&\le\sum_{\ell\in\mathcal E_n}D(\ell)
      \le e^{\varepsilon_n^{\mathrm{empty}}},\\
  0&\le\sum_{\ell\in\mathcal C_n}D(\ell)
      \le\varepsilon_n^{\mathrm{central}},\\
  0&\le\sum_{\ell\in\mathcal F_n}D(\ell)
      \le\varepsilon_n^{\mathrm{full}}.
\end{aligned}
\]
Every error sequence is deterministic, phase-independent, and tends to zero.
The lower bound in the first line is the term $D(0)=1$. Adding the disjoint
sums gives
\[
  1
  \le
  \sum_{0\le\ell_i\le k_i}D(\ell)
  \le
  e^{\varepsilon_n^{\mathrm{empty}}}
  +\varepsilon_n^{\mathrm{central}}
  +\varepsilon_n^{\mathrm{full}}
  =1+o(1),
\]
which proves (7.7) with one eventuality threshold for the complete phase.
